\documentclass[11pt]{article}
\usepackage{amsmath,amssymb,amsthm}
\usepackage[margin=1in]{geometry}

\newtheorem{theorem}{Theorem}
\newtheorem{lemma}[theorem]{Lemma}

\title{Problem 551: Sum of Digits Sequence}
\author{Project Euler Solutions}
\date{}

\begin{document}
\maketitle

\section{Problem Statement}

Let $a_0, a_1, a_2, \ldots$ be an integer sequence defined by $a_0 = 1$ and $a_n = a_{n-1} + s(a_{n-1})$ for $n \geq 1$, where $s(m)$ denotes the digit sum of~$m$ in base~10. Given that $a_{10^6} = 31054319$, find $a_{10^{15}}$.

\section{Mathematical Foundation}

\begin{theorem}[Digit-Sum Growth Bound]
For all $n \geq 0$, if $a_n$ has $d$ digits in base~10, then $a_{n+1} - a_n = s(a_n) \leq 9d$. Consequently, $a_n = O(n \log n)$.
\end{theorem}

\begin{proof}
Each base-10 digit is at most~9, and a $d$-digit number has exactly $d$ digits, so $s(m) \leq 9d$. Since $d = \lfloor \log_{10}(a_n) \rfloor + 1$, the increment per step is $O(\log a_n)$, yielding $a_n = O(n \log n)$ by summation.
\end{proof}

\begin{lemma}[Block Decomposition Principle]
Write $a_n$ in base~$10^B$, i.e., $a_n = \sum_{j=0}^{L} c_j \cdot 10^{jB}$ with $0 \leq c_j < 10^B$. Define $s_{\mathrm{upper}} = \sum_{j=1}^{L} s(c_j)$. Then the evolution of $c_0$ depends only on its current value and $s_{\mathrm{upper}}$, which remains constant until a carry propagates out of block~$c_0$.
\end{lemma}

\begin{proof}
The full digit sum is $s(a_n) = s(c_0) + s_{\mathrm{upper}}$. Adding $s(a_n)$ to $a_n$ modifies $c_0$ to $c_0 + s(c_0) + s_{\mathrm{upper}}$. If $c_0 + s(c_0) + s_{\mathrm{upper}} < 10^B$, no carry occurs and blocks $c_1, c_2, \ldots$ are unchanged, so $s_{\mathrm{upper}}$ is invariant.
\end{proof}

\begin{theorem}[Memoized Block Acceleration]
For block size~$B$, define $T(V, U) = (\Delta, V')$ where $V \in \{0, \ldots, 10^B - 1\}$ is the block value, $U$ is the upper digit sum, $\Delta$ is the number of steps until the first carry out of the block, and $V'$ is the new block value. Then $T$ is well-defined and computable by simulation of at most $10^B$ steps.
\end{theorem}

\begin{proof}
Since $a_n \geq 1$ for all $n$, each step increments $V$ by at least~1 (the digit sum of a positive integer is positive). Hence $V$ must reach~$10^B$ within at most~$10^B$ steps, producing a carry. The new value $V' = V_{\text{final}} - 10^B$ is uniquely determined.
\end{proof}

\section{Algorithm}

\begin{enumerate}
\item Represent $a_n$ as an array of 20 base-10 digits.
\item Maintain a memoization table indexed by (bottom-block value, upper digit sum).
\item At each iteration, look up or compute the block transition: how many steps until carry, and the resulting bottom-block value.
\item If the jump fits within the remaining step budget, apply it and propagate the carry into upper digits. Otherwise, fall back to single-step simulation.
\item Repeat until $10^{15}$ steps are completed.
\end{enumerate}

\section{Complexity Analysis}

\begin{itemize}
\item \textbf{Time:} The memoization table has $10^B \times U_{\max}$ entries. With $B = 4$ and $U_{\max} \leq 153$, this gives ${\sim}1.5 \times 10^6$ entries. Each is computed in $O(10^B)$ worst case. The main loop performs $O(N / \bar{\Delta})$ iterations where $\bar{\Delta}$ is the average jump size.
\item \textbf{Space:} $O(10^B \cdot U_{\max}) \approx O(1.5 \times 10^6)$.
\end{itemize}

\section{Answer}

\[
\boxed{73597483551591773}
\]

\end{document}
